\chapter{LÝ THUYẾT VỀ TINH CHỈNH VÀ CĂN CHỈNH MÔ HÌNH}
\label{chap:finetuning_alignment}

Trong chương~\ref{chap:transformer_llms} và chương~\ref{chap:pretraining}, chúng ta đã khám phá kiến trúc của các Mô hình Ngôn ngữ Lớn (LLMs) và cách chúng được \textbf{huấn luyện trước (pre-training)} trên một lượng dữ liệu văn bản khổng lồ. Quá trình này trang bị cho chúng một sự hiểu biết sâu sắc về ngôn ngữ, ngữ pháp, ngữ nghĩa và một lượng lớn kiến thức về thế giới. Tuy nhiên, một mô hình được huấn luyện trước giống như một sinh viên tốt nghiệp xuất sắc nhưng chưa có kinh nghiệm làm việc chuyên sâu: rất thông minh nhưng chưa phải là chuyên gia trong một lĩnh vực cụ thể nào.

Chương này sẽ tập trung vào giai đoạn thứ hai và cũng là giai đoạn quyết định để biến một mô hình nền tảng (foundation model) thành một ứng dụng hữu ích: giai đoạn \textbf{thích ứng (adaptation)}. Chúng ta sẽ tìm hiểu các kỹ thuật để “dạy chuyên sâu” cho các LLM, giúp chúng trở nên xuất sắc trong các tác vụ cụ thể, từ phân loại email đến việc tuân theo các chỉ dẫn phức tạp của con người.